% R61 — Registereintrag für Dok. 190 (Korrekturregister, append-only)
% Einzufügen als nächste Tabellenzeile nach R60.
% Autor: Johann Pascher, 23. Juli 2026

 R61 & 005 \S{}Baryon-Beispielrechnung (Proton) \newline (vgl.\ R50, P35) &
   Die gedruckte Proton-Rechnung
   $m_p = m_\mu \cdot K_\text{corr} \cdot QZ \cdot RG \cdot D_\text{baryon} \cdot f_\text{NN}$
   mit den Druckwerten $0{,}105658 \cdot 0{,}985 \cdot 0{,}532 \cdot 1{,}00007 \cdot 0{,}363 \cdot 1{,}00002$
   ergibt $0{,}0201$~GeV, nicht das angegebene Ergebnis $0{,}9381$~GeV
   (Fehlfaktor $46{,}672$). Die Formelzeile ist zudem dimensional offen
   ($m_\mu$ in GeV, $D_\text{baryon}$ in GeV $\Rightarrow$ GeV$^2$). &
   \textbf{Der Anker der Formelzeile ist $\pi^2/2$, nicht $m_\mu$.}
   \emph{(i) Diagnose.} Der Fehlfaktor identifiziert den Fehler als ein
   einziges Symbol: $(\pi^2/2)/m_\mu = 46{,}705$ trifft den Fehlfaktor
   auf $0{,}07\,\%$. \emph{(ii) Reparatur.} Mit
   $m_p = (\pi^2/2)\cdot K_\text{corr}\cdot QZ\cdot RG\cdot
   D_\text{baryon}\cdot f_\text{NN}$ --- $\pi^2/2$ dimensionslos,
   $D_\text{baryon}$ (enthält $0{,}5\,\Lambda_\text{QCD}$) alleiniger
   Einheitenträger --- schließt die Dimensionsbilanz, und die
   \emph{unveränderten} Druckwerte ergeben $0{,}93878$~GeV gegen
   PDG $0{,}93827$ ($\Delta = +0{,}05\,\%$) bzw.\ gegen das gedruckte
   Ergebnis $0{,}93810$ ($\Delta = +0{,}07\,\%$). Die Druckwerte waren
   mit dem $\pi^2/2$-Anker in sich konsistent; fehlerhaft war allein
   das Symbol $m_\mu$ in der Formelzeile. \emph{(iii) Strukturnotiz.}
   $\pi^2/2 = \text{Vol}(B^4)$ ist T$^4$-nativ;
   $(\pi^2/2)\cdot(16/\pi^2) = 8$ verbindet den Anker mit dem Kehrwert
   der D4-Packungsdichte (A270). Weitergehende Struktur (Aufspaltung
   $QZ\cdot K_\text{corr}$ unterbestimmt; geometrischer Kandidat
   $\pi/6$ für das Produkt; $\pi^3/12$-Verbindung zur Higgs-Formel
   über die Weyl-Zahl 192) wird eigenständig in A155
   \S{}Baryon-Kandidat geführt.
   \emph{(iv) Sekundärbefund.} Der Druckwert $K_\text{corr} = 0{,}985$
   passt nicht zur gedruckten Formel
   $K_\text{frak}^{D_f(1-(\xi/4)n_\text{eff})} = 0{,}9605$
   ($\Delta = 2{,}5\,\%$); die Beispielrechnung schließt nur mit
   $0{,}985$. Welcher Wert korrekt ist, entscheidet die ausstehende
   Herleitung. \emph{(v) Unberührt.} $f_\text{NN}$ ist ML-Kalibrierfaktor
   (numerisch $\approx 1$). Die Meson-Formel von Dok.~005 ist von
   diesem Eintrag nicht betroffen; eigenständige A-Serie-Nachfolger
   sind A155 (Meson und Baryon-Kandidat) und A270 (Sektor-Struktur).
   \emph{Prüfskript:} \texttt{python/Dok190\_Skripte/r61\_baryon\_anker.py}.
   \textbf{Die Quelldokumente werden nicht revidiert}
   (append-only; vgl.\ R50). \emph{Deklariert 23.~Juli 2026} \\
